% Solutions of 13-asymptotic-analysis.
\section*{Chapter \ref{ch:b1:asymptotic-analysis} --- Asymptotic Analysis and Recurrences}

\begin{solution}{exo:b1:asymptotic-analysis:1}
\omoutput[2-3]{code/bachelor-1/13-asymptotic-analysis/outputs/exo_count.txt}{The two counts.}
For each $i$ the inner loop runs $10 - i$ times: $10 + 9 + \dots + 1 = 55 = 10 \cdot 11/2$, in general $n(n + 1)/2$.
The second loop doubles $k$ until it reaches 1000: $k = 2^{10} = 1024$ after 10 steps, in general $\lceil \log_2 N
\rceil$ steps for a bound $N$ that is not a power of 2.
\end{solution}

\begin{solution}{exo:b1:asymptotic-analysis:2}
$(\log n)^{10} \ll \sqrt{n} \ll n^2 \log n \ll n^{2.5} \ll 1.01^n$. First, $(\log n)^{10}/\sqrt{n} = (\log n /
n^{1/20})^{10} \to 0$. Then $\sqrt{n}/(n^2 \log n) \to 0$ trivially. Next $n^2 \log n / n^{2.5} = \log n / n^{0.5} \to
0$. Last, $n^{2.5}/1.01^n \to 0$ since $1.01 > 1$.
\end{solution}

\begin{solution}{exo:b1:asymptotic-analysis:3}
$T(2) = 3$, $T(4) = 7$, $T(8) = 15$, $T(16) = 31$. Guess $T(n) = 2n - 1$: true for $n = 1$, and $2(2 \cdot \frac n2 - 1) + 1
= 2n - 1$. It is $\Theta(n)$, the \omterm{def:b1:induction-and-termination:rooted-tree}{leaf} case of the master theorem ($a = 2 > 2^0$).
\end{solution}

\begin{solution}{exo:b1:asymptotic-analysis:4}
(a) $a = 4 > 2^1$: leaves, $\Theta(n^2)$. (b) $a = 4 = 2^2$: even, $\Theta(n^2 \log n)$. (c) $a = 1 < 3^1$: root,
$\Theta(n)$. (d) $k = \frac12$ and $4^{1/2} = 2 = a$: even, $\Theta(\sqrt{n} \log n)$.
\end{solution}

\begin{solution}{exo:b1:asymptotic-analysis:5}
The statement $n^2 \in O(n)$ asks for \emph{one} constant $c$ valid for all large $n$. The ``induction'' proves $n^2 \leq
c_n n$ with a constant that grows at each step: $c_1 = 1$, $c_{n + 1} = c_n + 3$, so $c_n = 3n - 2$. What it proves is
$n^2 \leq (3n - 2)n$, which is true and says nothing. An induction for a $O$ statement must keep the same constant in
the hypothesis and the conclusion, and here it cannot: $n^2/n = n$ is unbounded.
\end{solution}

\begin{solution}{exo:b1:asymptotic-analysis:6}
The bound $O(n \log n)$ is right, but it is not a $\Theta$: the inner loop runs $\log_2 n$ times only for small $i$. For a
given $i$ it runs $\lceil \log_2 (n/i) \rceil$ times, and at least $t$ times exactly when $i \cdot 2^{t - 1} < n$, for
fewer than $n/2^{t - 1}$ values of $i$. Summing over $t \geq 1$, the count is below $n(1 + \frac12 + \frac14 + \cdots) =
2n$; and it is at least $n - 1$, since each $i < n$ runs once. So the \omterm{def:b1:asymptotic-analysis:cost}{cost} is $\Theta(n)$, and the output shows
$\text{count}/n$ approaching 2 while $\text{count}/(n \log_2 n)$ tends to 0.
\end{solution}

\begin{solution}{exo:b1:asymptotic-analysis:7}
\omcode{code/bachelor-1/13-asymptotic-analysis/exercises/nested_count.py}{4}{14}{Three nested loops}
Each iteration corresponds to one set $\{i, j, k\}$ of three distinct indices written in increasing order, and each such
set to one iteration: there are $\binom{n}{3} = n(n - 1)(n - 2)/6$ of them, in $\Theta(n^3)$. The test compares the
counter with the formula for $n < 12$.
\end{solution}

\begin{solution}{exo:b1:asymptotic-analysis:8}
Upper bound: each term is at most $n \log_2 n$, so the sum is at most $n^2 \log_2 n$. Lower bound: the terms with $i
\geq n/2$, at least $n/2$ of them, are at least $\frac n2 \log_2 \frac n2$, so the sum is at least $\frac{n^2}{4}(\log_2 n
- 1) \geq \frac{n^2}{8} \log_2 n$ for $n \geq 4$.
\end{solution}

\begin{solution}{exo:b1:asymptotic-analysis:9}
$T(n) = 1 + 2 + \dots + n = n(n + 1)/2$. For $U$, $U(n) + 1 = 2(U(n - 1) + 1)$, so $U(n) + 1 = 2^n$ and $U(n) = 2^n -
1$; both closed forms follow by induction, and the test checks them to $n = 59$. The master theorem covers subproblems
of size $n/b$, a constant fraction of $n$; here the size decreases by 1, and the \omterm{def:b1:asymptotic-analysis:recurrence}{recursion tree} has $n$ levels instead of
$\log_b n$.
\end{solution}

\begin{solution}{exo:b1:asymptotic-analysis:10}
With $q = a/b^k = 1$, the sum is $m = \log_b n$ and $a^m = n^{\log_b a} = n^k$, so $T(n) = n^k (T(1) + \log_b n)$ exactly.
For $n \geq b$, $\log_b n \geq 1$, so $n^k \log_b n \leq T(n) \leq (T(1) + 1)\, n^k \log_b n$: $T(n) \in \Theta(n^k \log
n)$, the base of the logarithm being irrelevant.
\end{solution}

\begin{solution}{exo:b1:asymptotic-analysis:11}
\emph{Increasing}, by strong induction: $T(2) = 2 > T(1)$; for $n \geq 2$, $\lfloor (n + 1)/2 \rfloor \geq \lfloor n/2
\rfloor$ and $\lceil (n + 1)/2 \rceil \geq \lceil n/2 \rceil$, all four at most $n$, so by the hypothesis $T(n + 1) \geq
T(n) + 1$. \emph{Order}: for $2^{m - 1} < n \leq 2^m$, $T(n) \leq T(2^m) = m 2^m \leq 2n(\log_2 n + 1)$ and $T(n) \geq
T(2^{m - 1}) = (m - 1)2^{m - 1} \geq \frac n2 (\log_2 n - 1)$, using the power-of-2 values $T(2^m) = m 2^m$ of
\cref{out:b1:asymptotic-analysis:recurrences}. So $T(n) \in \Theta(n \log n)$. \emph{Exact form}: the test checks $T(n)
= n\lceil \log_2 n \rceil + n - 2^{\lceil \log_2 n \rceil}$ for every $n < 3000$, and also $n \log_2 n \leq T(n) \leq n
\log_2 n + n$.
\end{solution}

\begin{solution}{exo:b1:asymptotic-analysis:12}
\omcode{code/bachelor-1/13-asymptotic-analysis/exercises/solve_rec.py}{7}{24}{Reading the case from the numbers}
The values $T(b^i)$ follow the recurrence directly, as integers. The ratio of the last two values is about $b^e$ for
$T \in \Theta(n^e)$, and about $b^k(1 + 1/m)$ in the even case: an exponent clearly above $k$ means the leaves dominate.
Otherwise the quotient $T(n)/n^k$ distinguishes the two other cases: it grows by about one per level in the even case
(it is $T(1) + \log_b n$) and converges in the root case. The test compares the answer with
\cref{lst:b1:asymptotic-analysis:master} on eleven recurrences.
\end{solution}

\begin{solution}{pb:b1:asymptotic-analysis:1}
\textbf{1.} For each $i$ the inner loop runs $n - 1 - i$ times: $(n - 1) + (n - 2) + \dots + 0 = n(n - 1)/2$, independent
of the values.

\textbf{2.} Every input of size $n$ \omterm{def:b1:asymptotic-analysis:cost}{costs} the same, so the three cases coincide, in $\Theta(n^2)$.

\textbf{3.} $25 \cdot 24/2 = 300$, $1225$, $4950$, $19\,900$ and $400 \cdot 399/2 = 79\,800$.

\textbf{4.} $\frac{2n(2n - 1)}{n(n - 1)} = \frac{2(2n - 1)}{n - 1} \to 4$; for $n = 25$ it is $1225/300 \approx 4.08$.

\textbf{5.} Return as soon as two equal values are found. Best case: the first two loans are of the same book, one
comparison; worst case: no duplicate, $n(n - 1)/2$ comparisons.

\textbf{6.} $\Theta(1)$: two indices and a counter.

\textbf{7.} In a sorted list the occurrences of a value form one block of consecutive positions; a value occurring at
least twice has two neighbours in its block, which the scan compares.

\textbf{8.} By strong induction, with $c = \lceil \log_2 n \rceil$: since $n \leq 2^c$, both halves have at most
$2^{c - 1}$ elements, so $T(n) \leq \lfloor \frac n2 \rfloor (c - 1) + \lceil \frac n2 \rceil (c - 1) + n - 1 < nc$.

\textbf{9.} At most $n \lceil \log_2 n \rceil$ comparisons to sort and $n - 1$ to scan: $O(n \log n)$.

\textbf{10.} $400 \log_2 400 \approx 3458$; the count is 3359, of which 399 for the scan: merge sort used about 2960, below
the bound, because a merge stops comparing when one half is exhausted.

\textbf{11.} $\frac{2n \log_2 (2n)}{n \log_2 n} = \frac{2(\log_2 n + 1)}{\log_2 n}$. For $n = 25$: $2(1 + 1/4.64) \approx
2.43$; for $n = 200$: $2(1 + 1/7.64) \approx 2.26$. The table shows 2.46 and 2.28.

\textbf{12.} The counts are not exactly $c\,n \log_2 n$: the merges depend on the data, and the scan adds a linear term.

\textbf{13.} One look-up per element of the list: $n$.

\textbf{14.} That a look-up and an insertion each \omterm{def:b1:asymptotic-analysis:cost}{cost} a bounded number of \omterm{def:b1:asymptotic-analysis:ram}{elementary operations}, which hash tables
achieve on average.

\textbf{15.} The double loop $\Theta(1)$; sorting $\Theta(n)$, for the copies made by merge sort; the set as many cells
as there are distinct values, at most $n$ (here at most 80).

\textbf{16.} \code{seen = [False] * 81}; for each \code{x}, test and set \code{seen[x]}: $n$ steps of constant \omterm{def:b1:asymptotic-analysis:cost}{cost},
plus 81 to create the list, $\Theta(n)$ time and $\Theta(1)$ space for a fixed range of identifiers.

\textbf{17.} $79\,800/400 \approx 200$ times fewer.

\textbf{18.} The model counts one dominant operation; each look-up \omterm{def:b1:asymptotic-analysis:cost}{costs} a bounded number of \omterm{def:b1:asymptotic-analysis:ram}{elementary operations}, so
the true \omterm{def:b1:asymptotic-analysis:cost}{cost} is within a constant factor of the count, which $\Theta$ ignores.

\textbf{19.} $c(2n)^k/(c\,n^k) = 2^k$.

\textbf{20.} Dividing numerator and denominator by $n^k$: $\frac{c\,2^k + d\,2^{k - 1}/n}{c + d/n} \to 2^k$.

\textbf{21.} Yes. $f(n) = n$ when $\lfloor \log_2 n \rfloor$ is even and $3n$ when it is odd is in $\Theta(n)$, but
$f(2n)/f(n)$ alternates between 6 and $2/3$.

\textbf{22.} \code{pairs}: 4.08, 4.04, 4.02, 4.01, quadratic; \code{sorting}: 2.46, 2.32, 2.31, 2.28, slightly above 2 and
decreasing, $n \log n$; \code{set}: exactly 2.00, linear.

\textbf{23.} The excess $2/\log_2 n$ decreases like a logarithm: it is still 0.2 at $n = 1024$ and 0.1 at about a
million.

\textbf{24.} \emph{The doubling ratio test}: when $n$ doubles, the ratio of the counts tends to 4 for $c\,n^2$ (question
19), to 2 for $c\,n$, and equals $2(1 + 1/\log_2 n)$ for $c\,n \log_2 n$ (question 11), which tends to 2 from above. On
the table the double loop gives 4.08 down to 4.01, sorting 2.46 down to 2.28, and the set 2.00: $\Theta(n^2)$, $\Theta(n
\log n)$ and $\Theta(n)$. The test is a diagnosis, not a proof: question 21 shows a $\Theta(n)$ \omterm{def:b1:asymptotic-analysis:cost}{cost} whose ratio never
settles.
\end{solution}
